\section{Storage: what survives}
\label{sec:storage}

Exactly one path spans both machines.

\begin{table}[!ht]
  \centering
  \renewcommand{\arraystretch}{1.25}
  \begin{tabularx}{\linewidth}{@{}l X X@{}}
    \toprule
    \textbf{Path} & \textbf{Shared with container?} &
    \textbf{Survives \cmd{gpu release}?} \\
    \midrule
    \cmd{/shared/\$USER} & yes, live and bidirectional & yes \\
    \cmd{\textasciitilde} on the developer VM & \textbf{no} & n/a \\
    \cmd{\textasciitilde} inside the container & it is \emph{camimo's} home & yes, on camimo \\
    container root filesystem & --- & no \\
    \bottomrule
  \end{tabularx}
  \caption{Which paths are shared and which survive a release.}
  \label{tab:gpu-storage}
\end{table}

This deserves emphasis because it is easy to get wrong. The claim command
bind-mounts \cmd{-v \$HOME:\$HOME}, but \cmd{\$HOME} is resolved \emph{on the
GPU host}. The home directory visible inside the container is camimo's
\cmd{/home/\$USER}, not the developer VM's. A file written to
\cmd{\textasciitilde} in one place is not visible in the other. Only
\cmd{-v /shared:/shared} is unconditional, and only \cmd{/shared} genuinely
spans both.

\medskip
\textbf{\cmd{/shared} is not a backup.} One unmirrored NVMe copied hourly to
another machine. Model weights belong there because they are large and
re-downloadable; anything whose loss would matter belongs in Git, pushed
elsewhere. It is also readable by the whole team, so it must never hold
credentials, keys, tokens, draft commercial terms, or material under NDA.
